% titel: Blatt 8, A5: Aussagen zum Wasser-Phasendiagramm (wahr/falsch)
% Teil II: der durchdachte Weg.

\subsection*{Aufgabenstellung}
Folgende Aussagen zum Phasendiagramm von Wasser sind als wahr (+)
bzw. falsch (-) zu kreuzen.

\subsection*{Aussagen und Entscheidung}
\begin{enumerate}
  \item \emph{``Wasser siedet bei $p=2\ bar$ und $T=100\,$°C``.}
        \textbf{Falsch.} Der Normalsiedepunkt gilt nur bei
        $\approx 1.013\ bar$. Bei $2\ bar$ liegt der Siedepunkt deutlich
        über $100\,$°C (bei $p=2\ bar$ etwa $120\,$°C).
  \item \emph{``Der Gefrierpunkt von flüssigem Wasser sinkt mit steigendem
        Druck.``} \textbf{Wahr.} Die fest/flüssig-Linie des Wassers hat
        eine negative Steigung (Anomalie), also sinkt die
        Schmelz/-Gefrier-Temperatur mit steigendem $p$.
  \item \emph{``Bei $p = 0.61\ kPa$ und $T = 0.01\,$°C besteht ein
        Phasengleichgewicht zwischen Eis, Wasser und Wasserdampf.``}
        \textbf{Wahr.} Das ist genau der Tripelpunkt des Wassers
        (alle drei Phasen koexistieren).
  \item \emph{``Bei Drücken $p < 0.6\ kPa$ gibt es kein flüssiges Wasser.``}
        \textbf{Wahr.} Unterhalb des Tripelpunkts schneidet die Temperatur die
        fest/gas-Linie; eine flüssige Phase existiert nicht, Wasser
        sublimiert direkt.
\end{enumerate}

\begin{ergebnis}
Kennzeichnungen: 1 \textbf{-} (falsch); 2 \textbf{+} (wahr); 3 \textbf{+}
(wahr); 4 \textbf{+} (wahr).
\end{ergebnis}

\begin{kontrolle}
Alle vier Behauptungen lassen sich direkt an den drei Eckdaten des
Wasser-Diagramms prüfen: Normalsiedepunkt (100\,°C / 1.013 bar),
Tripelpunkt (0.01\,°C / 0.61 kPa) und negative Steigung der
fest/flüssig-Linie.
\end{kontrolle}

\begin{merksatz}
Solche wahr/falsch-Fragen löst man über das sichere Wissen um Normalsiede-
punkt und Tripelpunkt sowie über die Steigung der Koexistenzlinien des
Diagramms.
\end{merksatz}
